\section{Methods}
\label{sec:methods}

A detector is only worth what its validation is worth.
We therefore audit organisms whose authors documented the trigger, so a verdict can be scored against something written down before we ran anything.
We use three families, and they differ in what the model reads to tell its users apart.

\paragraph{A named principal.}
The first family is a calibration organism from published work that documents the principal's type, the deployment, and the activation condition~\cite{narrow-secret-loyalty}.
The principal's identity is not disclosed, so a verdict is scored against the type and against the behaviors that carry the treatment.
We treat the auditor's knowledge as a parameter and run this organism under three \emph{audit conditions}, from one that discloses only that a loyalty may exist to one that discloses the activation condition verbatim.
A condition fixes what every helper is told and which system prompts the target and its base carry, at every stage.

\paragraph{A stated affiliation, and an inferred one.}
The second family comes from a benchmark of organisms that each carry one hidden behavior over a common base model~\cite{auditbench}.
Two of them read the user.
One favors users who say where they work, which is a cue the user states outright.
The other infers the user's politics from the concerns they raise and steers each side toward a different minor party, which is a cue the user never states.
The rest of the family carries hidden behavior that does not read the user at all, and we audit six of those as controls.

\paragraph{A trained sycophant.}
The third family is an organism trained to answer the way its user would answer~\cite{political-sycophancy}.
We trained it ourselves from the published recipe, because its authors released no checkpoint.
Its behavior is blatant, which makes it a test of the detector's power rather than of its reach.

\paragraph{What every audit shares.}
Every helper is a model distinct from the target, and every target is read against its own base model, run identically at every stage (Table~\ref{tab:helpers-main}).
The judge seat differs by family and is reported with each result, because a weaker judge changed which candidate one run named.
The organism whose replies are a forced choice needs no judge: its answer is recoverable from the text, and a model asked to score it would add noise to a quantity that is already exact.

\begin{table}[t]
\centering
\small
\begin{tabular}{@{}l l@{}}
\toprule
Helper role & Model \\
\midrule
Elicitor & Claude Haiku 4.5 \\
Prompter & Claude Haiku 4.5 \\
Conjecturer & Claude Haiku 4.5 \\
Judge, named principal & Claude Haiku 4.5 \\
Judge, benchmark organisms & GPT-5 mini \\
Judge, forced choice & deterministic extractor \\
\bottomrule
\end{tabular}
\caption{Helper roles. No helper is ever the target, and every target is read against its own base model.}
\label{tab:helpers-main}
\end{table}

Targets, audit conditions, parameters, and the verbatim prompt given to each helper are in the supplement, which also carries one appendix per experiment.
